% Báo cáo bài tập lớn môn Xử lý ảnh. Biên dịch bằng XeLaTeX (Overleaf: chọn XeLaTeX, file chính bao_cao.tex).
% Dòng lệnh: python3 bao_cao/latex/bien_dich.py   (cần xelatex hoặc tectonic; tự đồng bộ hình từ ket_qua/)
\documentclass[a4paper,12pt]{report}
\usepackage{fontspec}
\setmainfont{Times New Roman}
\usepackage[vietnamese]{babel}
\usepackage[top=2cm,bottom=2cm,left=3cm,right=2cm]{geometry}
\usepackage{amsmath,amssymb,graphicx,booktabs,array,setspace,hyperref,xurl,float}
\hypersetup{hidelinks,pdftitle={Nhận diện chữ số và hình dạng đơn giản bằng mạng neural}}
\onehalfspacing\raggedbottom
\setlength{\parskip}{6pt}
\graphicspath{{hinh/}}
\newcommand{\hinh}[3]{\begin{figure}[H]\centering\includegraphics[width=#2\textwidth]{#1}\caption{#3}\end{figure}}

\begin{document}

\begin{titlepage}
\centering
{\large HỌC VIỆN CÔNG NGHỆ BƯU CHÍNH VIỄN THÔNG}\\[1.2cm]
{\large BÁO CÁO BÀI TẬP LỚN}\\[0.4cm]
{\large Môn: Xử lý ảnh}\\[2cm]
{\Large\bfseries NHẬN DIỆN CHỮ SỐ VÀ HÌNH DẠNG ĐƠN GIẢN\\BẰNG MẠNG NEURAL}\\[3cm]
\begin{flushleft}\hspace{3cm}
\begin{tabular}{@{}ll}
Nhóm thực hiện: & \textit{[điền tên nhóm và thành viên]}\\
Giảng viên hướng dẫn: & \textit{[điền tên giảng viên]}\\
Thời điểm: & \textit{[điền tháng/năm]}
\end{tabular}
\end{flushleft}
\end{titlepage}

\tableofcontents
\clearpage

\chapter{Tổng quan}

\section{Đề tài và mục tiêu}
Đề tài xây dựng một hệ thống nhận dạng ảnh chứa \textbf{một chữ số viết tay} (0--9) hoặc \textbf{một hình học đơn giản}
(hình tròn, tam giác, hình chữ nhật hay vuông) bằng mạng neural tích chập (CNN). Mục tiêu gồm:
\begin{itemize}
\item cài đặt phần tiền xử lý ảnh bằng các thuật toán xử lý ảnh cơ bản của môn học (chuyển xám, lọc trung vị, phân ngưỡng Otsu,
hình thái học, thành phần liên thông, nội suy), đưa ảnh bất kỳ về đầu vào thống nhất cho mạng;
\item huấn luyện hai CNN nhỏ: một mạng cho chữ số (dữ liệu MNIST), một mạng cho hình học (dữ liệu tự sinh);
\item đóng gói thành ứng dụng web cho phép vẽ hoặc tải ảnh, xem kết quả cùng từng bước tiền xử lý;
\item đánh giá trung thực, nói rõ phần nào đã đo, phần nào chưa đo được.
\end{itemize}

\section{Phạm vi và giả định}
Mỗi ảnh chứa một đối tượng chính trên nền tương đối đồng đều. Người dùng chọn bài toán (chữ số hay hình học) trước khi nhận dạng; hai mạng
\emph{không} có lớp ``không phải'' nên ảnh sai loại vẫn bị gán một lớp. Toàn bộ thuật toán xử lý ảnh được tự cài bằng NumPy; thư viện Pillow chỉ
dùng để giải mã và ghi tệp ảnh, PyTorch dùng cho mạng neural.

\section{Kiến trúc hệ thống}
Ảnh đi qua một chuỗi tiền xử lý duy nhất, sau đó vào CNN tương ứng (Hình~\ref{hinh:luong}).

\begin{figure}[H]\centering
\begin{tabular}{c}
\fbox{\parbox{0.88\textwidth}{\centering Ảnh (vẽ trên khung / tải lên / tệp)}}\\[2pt]
$\downarrow$\\[2pt]
\fbox{\parbox{0.88\textwidth}{\centering Tiền xử lý: giới hạn kích thước $\rightarrow$ xám $\rightarrow$ (lọc trung vị) $\rightarrow$ Otsu $\rightarrow$ cực tính $\rightarrow$ thành phần liên thông lớn nhất\\ $\rightarrow$ hộp bao $\rightarrow$ thu nhỏ giữ tỉ lệ $\rightarrow$ căn giữa $28\times 28$ $\rightarrow$ kéo giãn độ sáng}}\\[2pt]
$\downarrow$\\[2pt]
\fbox{\parbox{0.42\textwidth}{\centering CNN chữ số (10 lớp)}}\ \ \fbox{\parbox{0.42\textwidth}{\centering CNN hình học (3 lớp)}}\\[2pt]
$\downarrow$\\[2pt]
\fbox{\parbox{0.88\textwidth}{\centering Lớp dự đoán và điểm softmax của mọi lớp}}
\end{tabular}
\caption{Luồng nhận dạng. Hai nhánh dùng chung tiền xử lý và chỉ khác lề căn đối tượng (20 px cho chữ số, 22 px cho hình).}
\label{hinh:luong}
\end{figure}

Mã nguồn chia theo chức năng thành các thư mục phẳng: \texttt{tien\_xu\_ly} (xử lý ảnh), \texttt{loi} (mạng, huấn luyện, đánh giá, dự đoán),
\texttt{bo\_du\_lieu} (dữ liệu) và \texttt{giao\_dien} (ứng dụng web); mỗi việc có một tệp chạy riêng ở thư mục gốc
(\texttt{huan\_luyen.py}, \texttt{nhan\_dang.py}, \texttt{danh\_gia.py}, \texttt{ung\_dung.py}). Chi tiết thiết kế xem tài liệu \texttt{bao\_cao/thiet\_ke.md}.

\section{Những điều báo cáo này \emph{không} khẳng định}
Các số liệu chính của hai CNN lấy từ MNIST (chữ viết tay có sẵn, đã chuẩn hóa) và từ dữ liệu hình học do chương trình sinh. Các thí nghiệm ở Chương~\ref{chuong:tichhop}
dùng ảnh do chương trình tạo để đo rủi ro, không phải độ chính xác trên ảnh thật.
Ảnh viết/vẽ tay thật của nhóm \textbf{chỉ có từ một người viết}: 58 ảnh chữ số và 43 ảnh hình, đã đủ số lượng tối thiểu đặt ra nhưng chưa có người viết khác, nên kết quả trên chúng
(Chương~\ref{chuong:tichhop}, mục~\ref{muc:thuanh}) chỉ là sơ bộ và không thể suy ra độ chính xác trên chữ viết tay nói chung.

\section{Bố cục}
Chương 2 trình bày tiền xử lý ảnh; Chương 3 và 4 trình bày hai CNN và kết quả; Chương 5 trình bày ứng dụng và đánh giá cả hệ thống;
Chương 6 kết luận và nêu hướng phát triển.

\chapter{Tiền xử lý ảnh}
\label{chuong:tienxuly}

\section{Mục tiêu và giả định}
Tiền xử lý đưa một ảnh bất kỳ chứa một đối tượng chính về dạng đầu vào thống nhất cho CNN: ảnh xám $28\times28$, giá trị thực trong $[0,1]$,
\textbf{nền bằng 0, nét sáng}, đối tượng được căn giữa và có kích thước cố định. Cách làm này loại bỏ các khác biệt không liên quan đến nhãn (màu sắc, độ sáng nền,
vị trí, kích thước, ảnh âm bản) để mạng chỉ cần học hình dạng. Giả định: một đối tượng, nền tương đối đồng đều, đối tượng không chiếm phần lớn viền ảnh.

\section{Các bước}

\subsection{Chuyển ảnh xám}
Ảnh màu RGB được chuyển theo độ chói $Y = 0{,}299R + 0{,}587G + 0{,}114B$. Ảnh có kênh alpha được ghép lên nền trắng trước khi chuyển xám:
$C' = \alpha C + (1-\alpha)\cdot 255$. Ảnh float được nhận diện thang đo $[0,1]$ hoặc $[0,255]$ từ giá trị lớn nhất.

\subsection{Giới hạn kích thước}
Thành phần liên thông được cài bằng duyệt BFS thuần Python nên rất chậm trên ảnh nhiều megapixel. Do đó ảnh có cạnh dài trên 512 px được thu nhỏ trước
bằng trung bình theo diện tích (mục~\ref{muc:dientich}); ảnh nhỏ hơn được giữ nguyên.

\subsection{Lọc trung vị (tùy chọn)}
Mỗi pixel được thay bằng trung vị của cửa sổ $k\times k$ ($k$ lẻ), biên được xử lý bằng cách lặp pixel ngoài cùng. Bộ lọc khử nhiễu muối tiêu rất tốt nhưng
\emph{xóa nét một pixel}, vì vậy \textbf{không bật mặc định}.

\subsection{Phân ngưỡng Otsu}
Với histogram 256 mức, ngưỡng $t$ được chọn để cực đại phương sai giữa hai lớp nền/nét \cite{otsu1979}:
\begin{equation}
t^{*}=\arg\max_{t}\ \omega_0(t)\,\omega_1(t)\,\bigl(\mu_0(t)-\mu_1(t)\bigr)^2 ,
\end{equation}
trong đó $\omega_i$ là tỉ lệ pixel và $\mu_i$ là mức xám trung bình của mỗi lớp. Otsu là ngưỡng toàn cục nên kém khi ánh sáng loang không đều.

\subsection{Cực tính}
Nét tối trên nền sáng và nét sáng trên nền tối được đưa về cùng một dạng. Chế độ tự động ước lượng nền từ trung vị mức xám của các pixel trên viền ảnh;
nếu tỉ lệ viền tối nằm trong $[0{,}35;\,0{,}65]$ chương trình cảnh báo vì quy tắc này có thể sai (ví dụ đối tượng chiếm viền). Người dùng có thể chỉ định cực tính.

\subsection{Hình thái học (tùy chọn)}
Co, giãn, mở, đóng với phần tử cấu trúc vuông $3\times3$. Co yêu cầu mọi vị trí của phần tử cấu trúc nằm trong đối tượng, giãn chỉ cần một vị trí chạm;
mở là co rồi giãn, đóng là giãn rồi co. Với mở/đóng, mặt nạ được đệm nền $r=\lfloor k/2\rfloor$ pixel trước khi ghép hai phép và chỉ cắt về kích thước gốc sau cả hai bước,
nhờ đó đối tượng chạm biên không bị co mất. Đây là một quy ước biên; quy ước ``cắt sau từng bước'' khác trong dải biên rộng $r$ và chúng được đối chiếu trong kiểm thử.
Hình thái học không bật mặc định vì có thể xóa nét hoặc làm dính nét.

\subsection{Thành phần liên thông lớn nhất}
Duyệt BFS theo 8 láng giềng để nét chéo vẫn là một thành phần, giữ thành phần nhiều pixel nhất và bỏ các chấm nhiễu rời. Hạn chế: nét bị đứt thành nhiều đoạn
chỉ còn đoạn lớn nhất (có tùy chọn tắt).

\subsection{Hộp bao, thu nhỏ và căn giữa}
Hộp bao nhỏ nhất của đối tượng được cắt ra và thu phóng với \emph{cùng một hệ số} cho hai chiều, $s = L/\max(H,W)$, để không làm méo tỉ lệ; $L$ là cạnh dài cho phép
($L=20$ cho chữ số, $L=22$ cho hình học trong canvas $28\times28$). Kết quả được đặt vào giữa canvas nền đen.

\subsection{Thu nhỏ theo diện tích}
\label{muc:dientich}
Khi \emph{phóng to} dùng nội suy song tuyến tính (giá trị tại tâm pixel đích được nội suy từ bốn pixel nguồn lân cận). Khi \emph{thu nhỏ mạnh} (ví dụ ảnh 400 px về 22 px) nội suy song tuyến tính chỉ
lấy mẫu vài điểm nên bỏ sót nét mảnh: nét biến thành các chấm rời và CNN nhận sai (xem Chương~\ref{chuong:tichhop}). Hệ thống vì vậy thu nhỏ bằng \textbf{trung bình theo diện tích}, tức lọc thông thấp
trước khi lấy mẫu: mỗi pixel đích là trung bình có trọng số của vùng nguồn mà nó phủ,
\begin{equation}
I'_{ij}=\sum_{u,v} W^{y}_{iu}\,I_{uv}\,W^{x}_{jv},\qquad W_{iu}=\frac{\operatorname{overlap}\bigl([i\,s,(i{+}1)s),[u,u{+}1)\bigr)}{s},
\end{equation}
với $s=n_{\text{vào}}/n_{\text{ra}}$ là bước thu nhỏ theo mỗi trục. Phép này bảo toàn giá trị trung bình và cho phép tỉ lệ không nguyên.

\subsection{Kéo giãn độ sáng}
Sau khi lấy trung bình, nét mảnh bị mờ (cường độ thấp), trong khi ảnh huấn luyện luôn có nét sáng nhất xấp xỉ 1. Ảnh được chia cho giá trị lớn nhất để nét sáng nhất bằng 1.

\section{Quy ước đầu vào của CNN}
\begin{table}[H]\centering
\begin{tabular}{lll}
\toprule
 & Chữ số & Hình học \\
\midrule
Kích thước, kiểu dữ liệu & $28\times28$, \texttt{float32} & giống nhau \\
Khoảng giá trị & $[0,1]$, nền 0, nét sáng & giống nhau \\
Cạnh dài tối đa của đối tượng & 20 px (lề 4) & 22 px (lề 3) \\
Căn giữa & theo tâm hộp bao & theo tâm hộp bao \\
\bottomrule
\end{tabular}
\caption{Quy ước đầu vào. Không trừ mean/chia std; không chia 255 lần thứ hai.}
\end{table}

\section{Kiểm thử}
Các hàm được đối chiếu với tham chiếu độc lập viết bằng vòng lặp hoặc đáp án dựng tay: chuyển xám (ba màu cơ bản cho 76/150/29; alpha 0/255/128 cho nền trắng/đen/xám 127),
lọc trung vị với duyệt láng giềng có sắp xếp (cửa sổ 1, 3, 5), Otsu với tính trực tiếp trung bình hai lớp cho cả 256 ngưỡng, bốn phép hình thái với vòng lặp Boolean (SE 1, 3, 5),
thành phần liên thông (đường chéo, hòa điểm, ảnh rỗng) và nội suy tại từng tọa độ. Ngoài ra có kiểm thử cho ảnh đảo màu, nền trong suốt, ảnh float, ảnh rất nhỏ, tham số sai
và ảnh không hợp lệ. Ở chế độ tương thích (song tuyến tính, không kéo giãn) đầu ra khớp từng pixel với đầu ra đã lưu của bản cài đặt đầu tiên (ví dụ chữ số có nhiễu, có trong bộ kiểm thử); khi gộp mã vào một package, chế độ này cũng được đối chiếu một lần với bản cũ trên 360 ca ngẫu nhiên, không có ca nào khác (phép đo một lần này không nằm trong bộ kiểm thử).

Đợt rà soát đầu tiên đã phát hiện và sửa ba lỗi: (1) mức xám của nét lấy từ ảnh gốc thay vì ảnh đã lọc nên nhiễu quay lại đầu ra; (2) vùng thêm vào nhờ phép đóng
không xuất hiện ở đầu ra; (3) phép đóng cắt phần giãn ngoài biên trước khi co nên làm mất vật chạm biên. Cả ba đều có kiểm thử hồi quy; ảnh trước/sau nằm trong
\texttt{ket\_qua/ra\_soat\_tien\_xu\_ly/}.

\section{Hạn chế}
Ngưỡng Otsu toàn cục kém trên nền loang sáng, bóng đổ mạnh hoặc nền hoa văn; tiền xử lý chỉ xử lý một đối tượng; cực tính tự động có thể sai khi đối tượng chiếm nhiều viền;
căn giữa theo hộp bao khác căn theo trọng tâm của MNIST (xem mục~\ref{muc:duongtienxuly}); chưa có tiêu chí tự động khẳng định ảnh đủ tốt cho CNN.

\chapter{CNN nhận dạng chữ số}
\label{chuong:chuso}

\section{Dữ liệu}
Dùng MNIST \cite{lecun1998}: 60\,000 ảnh huấn luyện và 10\,000 ảnh kiểm thử chính thức, ảnh xám $28\times28$, chữ số được căn theo trọng tâm. 6\,000 ảnh huấn luyện được tách làm
tập validation theo phân tầng (hạt 20261109), còn 54\,000 ảnh huấn luyện. 10\,000 ảnh test giữ nguyên và \textbf{không} dùng để chọn mô hình. Đã kiểm tra không có ảnh nào trùng pixel
giữa các tập. Ảnh được chia cho 255, không trừ mean/chia độ lệch chuẩn.

\hinh{chu_so_xem_truoc.png}{0.55}{Tám ảnh huấn luyện mỗi chữ số (MNIST, không phải ảnh ngoài).}

\section{Kiến trúc}
Hai khối Conv--ReLU--MaxPool rồi hai tầng kết nối đầy đủ:
$\text{Conv}(1\to16,3\times3)\to\text{ReLU}\to\text{MaxPool}(2)\to\text{Conv}(16\to32,3\times3)\to\text{ReLU}\to\text{MaxPool}(2)\to\text{Flatten}(800)\to\text{Dropout}(0{,}25)\to\text{Linear}(800\to64)\to\text{ReLU}\to\text{Linear}(64\to10)$,
tổng cộng \textbf{56\,714} tham số, nhỏ đủ để huấn luyện trên CPU. Mạng trả logits; softmax chỉ dùng khi dự đoán.

\section{Huấn luyện}
Hàm mất mát cross-entropy, bộ tối ưu Adam \cite{kingma2015} tốc độ học $10^{-3}$, giảm $0{,}3$ lần sau epoch 8 (StepLR), 12 epoch, batch 128, hạt 42, CPU.
Hai cấu hình giống hệt nhau (cùng kiến trúc, Dropout, hạt, ngân sách) và chỉ khác tăng cường dữ liệu:
\begin{itemize}
\item \textbf{cơ sở}: không tăng cường;
\item \textbf{cải tiến}: biến đổi affine ngẫu nhiên tự cài (xoay $\pm10^{\circ}$, dịch $\pm2$ px, co giãn $[0{,}9;1{,}1]$, nội suy song tuyến tính); không xoay $180^{\circ}$ hay lật vì có thể đổi $6\leftrightarrow9$.
\end{itemize}
Mỗi cấu hình lưu epoch có accuracy validation cao nhất (hòa thì loss validation thấp nhất). Cấu hình cuối được chọn \emph{chỉ bằng validation}, trước khi đọc test:
validation $98{,}88\%$ (cơ sở) và $99{,}02\%$ (cải tiến), nên chọn \textbf{cải tiến}.

\section{Kết quả}
\begin{table}[H]\centering
\begin{tabular}{lrrrr}
\toprule
Cấu hình & Đúng/10\,000 & Accuracy test & Khoảng Wilson 95\% & F1 macro \\
\midrule
Cơ sở & 9\,916 & $99{,}16\%$ & $[98{,}96;\,99{,}32]$ & $99{,}15\%$ \\
Cải tiến (được chọn) & 9\,915 & $99{,}15\%$ & $[98{,}95;\,99{,}31]$ & $99{,}14\%$ \\
\bottomrule
\end{tabular}
\caption{Kết quả trên 10\,000 ảnh test MNIST chính thức. Khoảng Wilson \cite{wilson1927} chỉ mô tả độ bất định do cỡ mẫu.}
\end{table}

Hai cấu hình chỉ khác nhau một ảnh. Kiểm định McNemar chính xác \cite{mcnemar1947} trên cùng 10\,000 ảnh: cơ sở sai/cải tiến đúng 37 ảnh, cơ sở đúng/cải tiến sai 38 ảnh, $p=1{,}0$.
Một hạt và một bộ đánh giá nên \textbf{không kết luận} tăng cường dữ liệu cải thiện độ chính xác trên MNIST chuẩn.

\hinh{chu_so_ma_tran.png}{0.6}{Ma trận nhầm lẫn của cấu hình được chọn trên tập test MNIST.}

Các cặp nhầm lẫn lớn nhất (lớp thật $\to$ lớp đoán): $9\to7$ (7 ảnh), $5\to3$ (6), $9\to4$ (6), $1\to2$ (4), $4\to9$ (4). Đây là các cặp có nét gần nhau, phù hợp với trực giác.

\hinh{chu_so_loi.png}{0.9}{24 trong 85 ảnh sai, softmax cao nhất trước (cấu hình được chọn).}

\section{Chẩn đoán độ bền}
Để thấy tác dụng của tăng cường dữ liệu, một tập con cố định 1\,000 ảnh test (100 ảnh mỗi chữ số) được biến đổi theo bốn điều kiện. Tập này lấy từ test nên
\textbf{không độc lập} với test và không dùng để chọn mô hình.
\begin{table}[H]\centering
\begin{tabular}{lrrrr}
\toprule
Cấu hình & Gốc & Dịch $\pm3$ px & Xoay $\pm20^{\circ}$ & Nhiễu Gauss $\sigma=0{,}15$ \\
\midrule
Cơ sở & $99{,}1\%$ & $89{,}6\%$ & $98{,}7\%$ & $99{,}3\%$ \\
Cải tiến & $99{,}0\%$ & $97{,}7\%$ & $98{,}3\%$ & $98{,}4\%$ \\
\bottomrule
\end{tabular}
\caption{Độ chính xác trên cùng 1\,000 ảnh test sau các biến đổi chẩn đoán.}
\end{table}
Tăng cường dữ liệu giúp rõ nhất ở phép dịch ($89{,}6\%\to97{,}7\%$), vì mô hình cơ sở chỉ thấy chữ số căn đúng trọng tâm; nó không tốt hơn ở mọi điều kiện.

\section{Hạn chế}
Kết quả ở chương này chỉ trên MNIST (một nguồn, một hạt). Trên 58 ảnh chữ số thật của một người viết, mô hình đạt $50/58$ (mục~\ref{muc:thuanh}), thấp hơn nhiều so với MNIST nhưng chỉ có một người viết nên chưa kết luận được; điểm softmax chưa hiệu chuẩn; mạng không từ chối chữ cái, hình học hay ảnh nhiều ký tự.
Đường vào qua tiền xử lý của hệ thống cho kết quả thấp hơn trên cùng ảnh (mục~\ref{muc:duongtienxuly}).

\chapter{CNN nhận dạng hình học}
\label{chuong:hinhhoc}

\section{Dữ liệu}
Không có bộ dữ liệu công khai phù hợp nên dữ liệu hình học được \textbf{sinh bằng chương trình}: bộ rasterizer giải tích bằng NumPy với siêu lấy mẫu $4\times4$, mỗi ảnh gốc $28\times28$ chứa một hình có
độ dày nét 1--3 px, góc quay $0^{\circ}$--$360^{\circ}$, đường kính ngoại tiếp 7--21 px, nhiễu Gauss nhẹ và nhiễu muối tiêu. Ba lớp: \emph{tròn}, \emph{tam giác}, \emph{chữ nhật hoặc vuông}.

Để mạng không chỉ học các hình quá lý tưởng, bộ sinh được làm đa dạng:
\begin{table}[H]\centering
\begin{tabular}{lp{0.66\textwidth}}
\toprule
Lớp & Các dạng được sinh \\
\midrule
Tam giác & đều 25\%, vuông 20\%, cân 20\%, bất kỳ 35\%; mọi góc trong $\ge18^{\circ}$ \\
Chữ nhật & tỉ lệ cạnh $U(0{,}3;1)$, 25\% là hình vuông; cạnh ngắn $\ge$ độ dày nét $+2$ px để lõi không bị lấp đặc \\
Tròn & 50\% elip (tỉ lệ trục $U(0{,}7;1)$), 50\% lượn sóng bán kính tới 5\%, phần còn lại tròn hoàn hảo \\
Rung tay & 50\% tam giác/chữ nhật có nhiễu Gauss lên đỉnh, $\sigma$ tới 2,5\% kích thước \\
\bottomrule
\end{tabular}
\caption{Độ đa dạng của bộ sinh dữ liệu hình học.}
\end{table}

Tổng cộng 13\,500 ảnh chia phân tầng 70/15/15 thành 9\,450 train, 2\,025 validation, 2\,025 test (mỗi lớp 675 ảnh test), các tập không giao nhau và không có ảnh trùng.
Ngoài ra có bộ \emph{stress} 300 ảnh khó hơn (nhiễu mạnh, mờ, có cột bị mất, đỉnh rung 3{,}5\%, hình tròn luôn méo) chỉ dùng để đo độ bền, không dùng huấn luyện hay chọn mô hình.
Mỗi ảnh gốc được đưa qua \emph{chính} tiền xử lý của hệ thống (Chương~\ref{chuong:tienxuly}) để tạo đầu vào CNN, nên huấn luyện và nhận dạng dùng chung một đường xử lý.

\hinh{hinh_hoc_xem_truoc.png}{0.75}{Ảnh hình học tổng hợp (gốc $28\times28$).}

\section{Kiến trúc và huấn luyện}
Cùng bộ khung với mạng chữ số nhưng đầu ra 3 lớp. Cấu hình \emph{cơ sở}: không BatchNorm, không Dropout, không tăng cường dữ liệu (56\,259 tham số). Cấu hình \emph{cải tiến}: thêm BatchNorm \cite{ioffe2015} sau hai
tầng tích chập, Dropout $0{,}3$ \cite{srivastava2014} và tăng cường dữ liệu (xoay $\pm20^{\circ}$, dịch $\pm1{,}5$ px, co giãn, đổi độ sáng, nhiễu Gauss $0{,}012$) (56\,355 tham số).
Adam $10^{-3}$, batch 64, 20 epoch, hạt 42, CPU; chọn epoch và cấu hình cuối bằng validation như ở Chương~\ref{chuong:chuso}.
Cấu hình cải tiến thay đồng thời ba yếu tố (BatchNorm, Dropout, tăng cường) nên không quy được khác biệt cho riêng yếu tố nào.

\section{Kết quả}
\begin{table}[H]\centering
\begin{tabular}{lrrrr}
\toprule
Cấu hình & Validation & Test (2\,025 ảnh) & F1 macro & Stress (300 ảnh) \\
\midrule
Cơ sở & $99{,}75\%$ & $2014/2025=99{,}46\%$ & $99{,}46\%$ & $258/300=86{,}0\%$ \\
Cải tiến (được chọn) & $99{,}90\%$ & $2019/2025=99{,}70\%$ & $99{,}70\%$ & $270/300=90{,}0\%$ \\
\bottomrule
\end{tabular}
\caption{Kết quả hình học; cấu hình cuối chọn bằng validation. Test và stress là dữ liệu tổng hợp.}
\end{table}
Khoảng Wilson 95\% của accuracy test: cơ sở $[99{,}03;\,99{,}70]\%$, cải tiến $[99{,}36;\,99{,}86]\%$. McNemar chính xác ghép cặp: cơ sở sai/cải tiến đúng 6 ảnh, ngược lại 1 ảnh, $p=0{,}125$
(chưa đủ để khẳng định khác biệt). Theo lớp, cấu hình được chọn có precision/recall từ $99{,}56\%$ đến $100\%$.
Sáu ảnh test sai gồm hai elip, một elip lượn sóng, hai chữ nhật và một hình vuông; các nhầm lẫn là tròn$\to$chữ nhật (3) và chữ nhật$\to$tam giác (3).

\hinh{hinh_hoc_ma_tran.png}{0.5}{Ma trận nhầm lẫn của cấu hình được chọn trên tập test hình học.}

Trên bộ stress, các nhầm lẫn lớn nhất là chữ nhật$\to$tam giác (13 ảnh) và tròn$\to$tam giác (9), tròn$\to$chữ nhật (7). Xem lại 24 lỗi có softmax cao nhất cho thấy phần lớn là hình \emph{bị đứt nét} (bộ stress xóa một cột pixel): bước giữ thành phần liên thông lớn nhất chỉ còn một cung hoặc một góc, và mạng đoán sai trên phần còn lại. Đây là hệ quả của giới hạn của bước liên thông chứ không chỉ của mạng.

\section{Độ bền theo góc xoay và độ dày nét}
Quét ghép cặp 40 đối tượng gốc $\times$ 12 góc $\times$ 3 độ dày nét (tam giác và chữ nhật của bộ stress, 480 ảnh mỗi độ dày, mỗi mô hình), thực hiện sau khi đã chọn mô hình:
\begin{table}[H]\centering
\begin{tabular}{lrrr}
\toprule
Độ dày nét (px) & 1 & 2 & 3 \\
\midrule
Cơ sở & $90{,}8\%$ & $92{,}1\%$ & $89{,}0\%$ \\
Cải tiến & $94{,}2\%$ & $94{,}8\%$ & $95{,}0\%$ \\
\bottomrule
\end{tabular}
\caption{Độ chính xác trong quét ghép cặp (dữ liệu tổng hợp, không phải hình người vẽ).}
\end{table}

\hinh{hinh_hoc_quet.png}{0.95}{Độ chính xác theo góc xoay và độ dày nét (mỗi ô 40 đối tượng khớp nhau).}

\section{So sánh dữ liệu huấn luyện}
\label{muc:sosanhdulieu}
Để kiểm tra việc làm đa dạng bộ sinh có thật sự cần thiết, cấu hình cải tiến được huấn luyện thêm trên hai bộ dữ liệu khác với cùng kiến trúc, công thức, hạt và tiền xử lý:
\emph{đơn giản} (tam giác đều, chữ nhật tỉ lệ 0,55--1, hình tròn hoàn hảo, không rung đỉnh; 4\,500 ảnh) và \emph{đa dạng} với cùng 4\,500 ảnh. Cả ba mô hình (kèm mô hình chính 13\,500 ảnh) được đo trên các bộ giống nhau.

\begin{table}[H]\centering\small
\begin{tabular}{lrrrr}
\toprule
Dữ liệu huấn luyện & Test đơn giản & Test đa dạng & Stress đa dạng & Hình vẽ Pillow (TB) \\
\midrule
Đơn giản, 4.500 ảnh & 100,0\% & 96,9\% & 80,3\% & 91,5\% \\
Đa dạng, 4.500 ảnh & 99,9\% & 99,5\% & 85,3\% & 99,7\% \\
Đa dạng, 13.500 ảnh (mô hình chính) & 100,0\% & 99,7\% & 90,0\% & 99,6\% \\
\bottomrule
\end{tabular}
\caption{Ảnh hưởng của dữ liệu huấn luyện (tất cả là dữ liệu tổng hợp).}
\end{table}

\begin{table}[H]\centering\small
\begin{tabular}{lrrr}
\toprule
Họ hình (Pillow) & Đơn giản 4\,500 & Đa dạng 4\,500 & Đa dạng 13\,500 \\
\midrule
Tròn & 100,0\% & 100,0\% & 100,0\% \\
Tam giác đều & 100,0\% & 100,0\% & 100,0\% \\
Tam giác vuông & 100,0\% & 100,0\% & 99,4\% \\
Tam giác bất kỳ & 100,0\% & 98,1\% & 99,4\% \\
Chữ nhật (0,55--1) & 100,0\% & 100,0\% & 100,0\% \\
Chữ nhật dài (0,25--0,45) & 48,8\% & 100,0\% & 98,8\% \\
\bottomrule
\end{tabular}
\caption{Độ chính xác theo họ hình trên ảnh vẽ Pillow (trung bình 4 độ dày nét, 160 ảnh mỗi họ).}
\end{table}
\hinh{so_sanh_du_lieu_hinh.png}{0.85}{So sánh dữ liệu huấn luyện hình học.}

Bộ test của chính dữ liệu đơn giản đã bão hòa ($100\%$) nên không phân biệt được mô hình tổng quát với mô hình học thuộc phân bố hẹp; trên test đa dạng mô hình đó chỉ đạt $96{,}9\%$.
Dữ liệu đa dạng cải thiện rõ nhất ở \textbf{chữ nhật dài} ($48{,}8\%\to100\%$ với 4\,500 ảnh, $98{,}8\%$ với 13\,500 ảnh); các họ còn lại gần như không đổi, vì tam giác vuông và tam giác bất kỳ
đã được nhận đúng ngay cả khi mô hình chỉ học tam giác đều (với tiền xử lý hiện tại). Gấp ba số ảnh tăng độ chính xác trên bộ stress ($85{,}3\%\to90{,}0\%$) nhưng không thấy khác biệt trên hình vẽ Pillow ($99{,}7\%$ so với $99{,}6\%$, trong sai số của 960 ảnh).
Một hạt cho mỗi cấu hình, nên các chênh lệch nhỏ chưa đủ ý nghĩa thống kê.

\section{Hạn chế}
Mọi số liệu ở chương này là dữ liệu tổng hợp cùng họ bộ sinh; bộ stress và bộ quét cũng do cùng bộ sinh tạo ra nên accuracy cao \textbf{không bảo đảm} hiệu quả trên hình người vẽ.
Hình tròn méo nặng kèm nhiễu mạnh và chữ nhật quá dài vẫn dễ sai: trong phép thử chữ nhật nét 3--8 px, tỉ lệ cạnh 0,22 và 0,31 nhận đúng, từ 0,15 trở xuống (ngoài dải huấn luyện 0,3--1) bị đoán thành tam giác. Mạng chỉ có ba lớp và không từ chối ngôi sao, hình thang hay chữ cái.

\chapter{Ứng dụng và đánh giá cả hệ thống}
\label{chuong:tichhop}

\section{Ứng dụng web}
Ứng dụng chạy cục bộ (máy chủ chỉ lắng nghe địa chỉ \texttt{127.0.0.1}, không gửi ảnh ra ngoài) bằng lệnh \texttt{python3 ung\_dung.py}. Người dùng chọn bài toán, vẽ trên khung hoặc tải, kéo thả,
dán ảnh; ứng dụng hiển thị lớp dự đoán, điểm softmax của mọi lớp và ảnh sau từng bước tiền xử lý (xám, lọc, Otsu, mặt nạ, cắt, thu nhỏ, đầu vào CNN $28\times28$), kèm các tùy chọn cực tính, lọc trung vị
và chỉ giữ thành phần lớn nhất. Ảnh lỗi (trống, không tách được đối tượng, nét quá mảnh) trả về thông báo giải thích thay vì kết quả. Ảnh chụp lớn được thu nhỏ trước khi xử lý.
Kiểm thử đã chạy máy chủ thật trên cổng tạm cho các luồng nhận dạng, tùy chọn, yêu cầu sai và thu ảnh.

\section{Đường qua tiền xử lý thống nhất với ảnh MNIST}
\label{muc:duongtienxuly}
CNN chữ số được huấn luyện trên ảnh MNIST gốc, vốn căn theo trọng tâm. Ứng dụng lại cắt theo hộp bao, thu phóng về 20 px và căn theo tâm hộp bao. Trên cùng 1\,000 ảnh test MNIST (100 ảnh mỗi chữ số),
đưa trực tiếp vào CNN đạt $99{,}0\%$, qua toàn bộ tiền xử lý thống nhất đạt $97{,}6\%$, giảm $1{,}4$ điểm phần trăm. Phép đo này thay đổi nhiều bước cùng lúc, nên chưa thể quy mức giảm cho riêng cách căn giữa;
muốn biết phần do cách căn cần so sánh riêng căn theo tâm hộp bao và căn theo trọng tâm trên cùng cấu hình (mục~\ref{muc:huongphattrien}).

\section{Chữ số và hình vẽ nét mảnh}
\label{muc:netmanh}
Khi người dùng vẽ bằng bút nhỏ trên khung vẽ lớn, nét chỉ bằng vài phần trăm kích thước đối tượng. Sau khi thu nhỏ về 20--22 px nét còn dưới một pixel; nếu thu nhỏ bằng nội suy song tuyến tính
nét bị rách và CNN nhận sai. Hai thí nghiệm đo hiện tượng này với ba cấu hình khâu thu nhỏ và chuẩn hóa: \emph{song tuyến tính}, \emph{trung bình theo diện tích} và \emph{diện tích kèm kéo giãn độ sáng}
(cấu hình mặc định). Cả hai dùng ảnh \textbf{do chương trình tạo}, vì vậy chỉ cho biết mức rủi ro.

\subsection{Chữ số}
150 ảnh test MNIST được phóng to, làm mỏng nét về bộ xương một pixel bằng thuật toán Zhang--Suen \cite{zhang1984} rồi giãn trở lại đến độ dày mong muốn, đặt ngẫu nhiên vào khung 300--480 px.

\begin{table}[H]\centering
\begin{tabular}{lrrr}
\toprule
Độ dày nét / chiều cao chữ số & Song tuyến tính & Diện tích & Diện tích + kéo giãn \\
\midrule
12\% & 99,3\% & 99,3\% & 99,3\% \\
8\% & 98,7\% & 98,7\% & 98,7\% \\
5\% & 98,0\% & 98,7\% & 98,7\% \\
3\% & 90,7\% & 98,0\% & 98,0\% \\
1,5\% & 52,7\% & 82,0\% & 95,3\% \\
\bottomrule
\end{tabular}
\caption{Độ chính xác chữ số theo độ dày nét (150 ảnh mỗi dòng; ảnh tổng hợp từ MNIST).}
\end{table}
\hinh{net_manh_chu_so.png}{0.7}{Chữ số nét mảnh: ba cấu hình thu nhỏ/chuẩn hóa.}

Với nét dày (từ 8\% trở lên) ba cấu hình như nhau, nên cách mới không làm kém đi trường hợp bình thường. Khi nét chỉ còn 3\% chiều cao, song tuyến tính giảm xuống $90{,}7\%$; ở $1{,}5\%$ chỉ còn $52{,}7\%$.
Trung bình diện tích giữ được nét nhưng làm nó mờ, nên cần thêm bước kéo giãn độ sáng ($82{,}0\%\to95{,}3\%$ ở $1{,}5\%$).

\subsection{Hình học}
Sáu họ hình (tròn, tam giác đều, vuông, bất kỳ, chữ nhật tỉ lệ vừa và dài) được vẽ bằng Pillow kiểu Paint, có rung tay, kích thước ảnh 300--500 px và bốn độ dày nét
(1,2\%, 2,5\%, 5\%, 9\% kích thước hình), 40 ảnh mỗi ô. Mã vẽ độc lập với bộ sinh dữ liệu huấn luyện nhưng \textbf{cùng các họ hình} mà mô hình đã được dạy.

\begin{table}[H]\centering\small
\begin{tabular}{lrrrrrr}
\toprule
 & \multicolumn{3}{c}{Nét mảnh nhất (1,2\%)} & \multicolumn{3}{c}{Trung bình 4 độ dày} \\
\cmidrule(lr){2-4}\cmidrule(lr){5-7}
Họ hình & S.t.t. & Diện tích & + kéo giãn & S.t.t. & Diện tích & + kéo giãn \\
\midrule
Tròn & 5\% & 100\% & 100\% & 66,9\% & 100,0\% & 100,0\% \\
Tam giác đều & 85\% & 100\% & 100\% & 95,6\% & 100,0\% & 100,0\% \\
Tam giác vuông & 85\% & 100\% & 100\% & 88,8\% & 100,0\% & 99,4\% \\
Tam giác bất kỳ & 85\% & 100\% & 100\% & 95,0\% & 99,4\% & 99,4\% \\
Chữ nhật (tỉ lệ 0,55--1) & 90\% & 60\% & 100\% & 97,5\% & 90,0\% & 100,0\% \\
Chữ nhật dài (0,25--0,45) & 38\% & 5\% & 98\% & 80,0\% & 75,6\% & 98,8\% \\
\midrule
Trung bình toàn bộ & \multicolumn{3}{c}{} & 87,3\% & 94,2\% & 99,6\% \\
\bottomrule
\end{tabular}
\caption{Hình vẽ Pillow có rung tay (S.t.t.: nội suy song tuyến tính). Mỗi ô $40$ ảnh; ảnh tổng hợp, không phải hình người vẽ thật.}
\end{table}
\hinh{ve_pil_hinh_hoc.png}{0.85}{Hình vẽ Pillow có rung tay: độ chính xác trung bình theo họ hình.}

Hình tròn nét mảnh nhất bị nhận sai gần như hoàn toàn ($5\%$) với song tuyến tính vì vòng tròn bị rách thành các chấm. Trung bình diện tích khắc phục điều này nhưng, nếu không kéo giãn độ sáng,
làm hỏng chữ nhật nét mảnh ($60\%$ và $5\%$) vì nét sau khi lấy trung bình quá mờ so với ảnh huấn luyện. Chỉ khi kết hợp cả hai bước mới đạt $99{,}6\%$ trung bình.
Hai bước này vì vậy được dùng làm mặc định.

\section{Thu ảnh thật và quy trình đánh giá}
\label{muc:thuanh}
Ứng dụng có sẵn trang thu ảnh (\texttt{python3 ung\_dung.py -{}-thu-anh}): người thu chọn \textbf{nhãn thật trước}, vẽ trên khung hoặc tải ảnh, đặt mã người viết tự đặt. Trang \emph{không hiển thị dự đoán}
để tránh sửa ảnh hay nhãn theo mô hình. Mỗi ảnh được lưu kèm tệp mô tả (nhãn, nguồn, thời điểm, SHA-256); ảnh trống hoặc không tách được đối tượng bị từ chối lúc thu. Lệnh
\texttt{python3 danh\_gia.py -{}-bai-toan ca\_hai} đánh giá trên các ảnh này bằng \emph{đúng đường xử lý của ứng dụng}, bỏ qua ảnh trùng byte, tính ảnh không xử lý được là sai và đếm riêng.
Yêu cầu tối thiểu để kết quả được coi là đủ: chữ số 50 ảnh (ít nhất 5 mỗi chữ số), hình học 30 ảnh (ít nhất 10 mỗi hình), nên có nhiều người viết/vẽ. Các ảnh này chỉ dùng để \emph{đánh giá},
không dùng để huấn luyện hay chỉnh mô hình đã chốt.

\subsection{Kết quả trên ảnh thật (một người viết)}
Đã thu 58 ảnh chữ số và 43 ảnh hình của \textbf{một người}: viết/vẽ bằng bút mực trên giấy, chụp ảnh các tờ giấy, rồi cắt từng ký hiệu thành một ảnh riêng
(chừa lề quanh nét); có ba tờ giấy. Ảnh chụp có nét chữ mặt sau tờ giấy lộ mờ qua giấy và ánh sáng không đều. Nhãn được ghi theo thứ tự đọc \emph{trước} khi chạy dự đoán; sau đó người viết đã rà soát lại toàn bộ nhãn.
Bộ này \textbf{đã đạt ngưỡng số lượng tối thiểu} ở trên (trạng thái \texttt{ready}, chỉ xét số lượng): chữ số 58 ảnh, mỗi chữ số ít nhất 5 ảnh; hình 43 ảnh (16 tròn, 17 tam giác, 10 chữ nhật/vuông).
\emph{Nhưng chỉ có một người viết}, nên kết quả không đại diện cho chữ viết tay nói chung.

\begin{table}[H]\centering
\begin{tabular}{lrrl}
\toprule
Bài toán & Cơ sở & Cải tiến (được chọn) & Khoảng Wilson 95\% (cải tiến) \\
\midrule
Chữ số ($58$ ảnh) & $49/58$ & $50/58$ ($86{,}2\%$) & $[75{,}1;\,92{,}8]\%$ \\
Hình học ($43$ ảnh) & $42/43$ & $43/43$ ($100\%$) & $[91{,}8;\,100]\%$ \\
\bottomrule
\end{tabular}
\caption{Ảnh thật của một người viết; khoảng Wilson chỉ mô tả độ bất định do cỡ mẫu, không nói về khác biệt giữa người viết.}
\end{table}

Với hình học, cả 43 ảnh được nhận đúng, gồm tam giác nhiều hướng (đỉnh lên, đỉnh xuống, nằm ngang), hình tròn lớn và nhỏ, chữ nhật rất dài và 3 hình tô đặc (hệ thống chỉ được huấn luyện trên hình đường viền);
mô hình cơ sở sai một ảnh (hình vuông tô đặc bị nhận là tam giác). Mẫu chữ nhật/vuông chỉ có 10 ảnh.
Với chữ số, $50/58$ thấp hơn nhiều so với $99{,}15\%$ trên MNIST. Tám ảnh sai là $2\to1$, $4\to2$, $4\to9$, $7\to2$, $8\to3$ và ba ảnh $9\to3$; chữ số 9 chỉ đúng $3/6$, chữ số 4 đúng $3/5$.
Theo từng tờ giấy kết quả là $25/30$, $20/23$ và $5/5$ (tờ cuối chỉ có chữ số 3). Quan sát ảnh đầu vào CNN bằng mắt (chưa phân tích định lượng):
cả ba chữ số 9 sai được viết kiểu chữ ``g'' có đuôi móc; chữ số 8 viết hẹp; một chữ số 4 viết đầu kín trông như số 9; một chữ số 7 có gạch ngang giữa bị nhận là 2 với độ tin cậy softmax thấp; ở một chữ số 4 khác,
nét sổ đứng không dính vào nét chéo nên bị bước giữ thành phần liên thông lớn nhất bỏ mất.
Tắt bước đó trên 30 ảnh của tờ thứ nhất làm số 4 đúng nhưng sai thêm hai ảnh khác ($24/30$), nên đây không phải lỗi sửa được bằng một tùy chọn.
Kết quả cho thấy khoảng cách giữa MNIST và chữ viết tay thật đáng được khảo sát tiếp (kiểu viết, cách căn chữ, ảnh chụp), chưa cho phép kết luận về nguyên nhân.
Các ảnh này chỉ dùng để đánh giá, không dùng để huấn luyện hay chỉnh mô hình.

\section{Hạn chế}
Ngưỡng Otsu toàn cục kém khi nền loang sáng nặng; chọn chữ số hay hình học là thao tác tay và hai mạng không từ chối ảnh sai loại; các thí nghiệm nét mảnh ở mục~\ref{muc:netmanh} dùng ảnh tổng hợp, một hạt và số ảnh vừa phải
(150 ảnh chữ số mỗi mức độ dày, 40 ảnh mỗi ô hình học), nên các chênh lệch vài điểm phần trăm không nên được diễn giải quá mức; mục~\ref{muc:thuanh} là phép đánh giá riêng trên ảnh viết/vẽ thật nhưng mới có một người.

\chapter{Kết luận}

\section{Kết quả đạt được}
Đề tài đã xây dựng một hệ thống nhận dạng chữ số viết tay và hình học đơn giản hoàn chỉnh từ ảnh đầu vào đến kết quả:
\begin{itemize}
\item phần tiền xử lý ảnh tự cài bằng NumPy với các thuật toán của môn học (xám, lọc trung vị, Otsu, hình thái học, thành phần liên thông, nội suy), được kiểm thử bằng tham chiếu độc lập;
\item CNN chữ số 56\,714 tham số đạt $9915/10\,000$ ($99{,}15\%$) trên tập test MNIST chính thức, cấu hình được chọn bằng validation;
\item CNN hình học 56\,355 tham số, huấn luyện trên 13\,500 ảnh tổng hợp đa dạng, đạt $2019/2025$ ($99{,}70\%$) trên tập test tổng hợp và $270/300$ ($90{,}0\%$) trên bộ stress;
\item ứng dụng web cho phép vẽ hoặc tải ảnh, xem kết quả và từng bước tiền xử lý, cùng công cụ thu ảnh thật để đánh giá.
\end{itemize}
Phát hiện đáng chú ý nhất của thực nghiệm là \textbf{khâu thu nhỏ ảnh}: nội suy song tuyến tính làm nét mảnh bị rách và CNN nhận sai (hình tròn nét mảnh gần như luôn sai; chữ số nét 1,5\% chiều cao chỉ còn $52{,}7\%$).
Thu nhỏ theo diện tích kèm kéo giãn độ sáng đưa độ chính xác trung bình trên hình vẽ có rung tay từ $87{,}3\%$ lên $99{,}6\%$ mà không làm kém ảnh nét dày.

\section{Những gì chưa làm được}
\begin{itemize}
\item \textbf{Ảnh viết/vẽ tay thật chỉ có từ một người viết}: 58 ảnh chữ số và 43 ảnh hình, đã đủ số lượng tối thiểu nhưng chưa có người viết khác. Trên bộ này chữ số đạt $50/58$ ($86{,}2\%$, khoảng Wilson $[75{,}1;\,92{,}8]\%$), thấp hơn nhiều so với $99{,}15\%$ trên MNIST; hình học đạt $43/43$ (khoảng Wilson $[91{,}8;\,100]\%$). Chỉ có một người viết nên chưa suy ra được độ chính xác chung. Các con số còn lại là MNIST và dữ liệu tổng hợp; các phép thử nét mảnh và hình vẽ Pillow dùng ảnh do chương trình tạo.
\item Chọn bài toán là thao tác tay; hai mạng không có lớp ``không phải'' nên ảnh sai loại vẫn bị gán một lớp rất tự tin; softmax chưa hiệu chuẩn.
\item Mỗi cấu hình chỉ huấn luyện với một hạt; các chênh lệch giữa cơ sở và cải tiến chưa đủ ý nghĩa thống kê.
\item Phép phân ngưỡng toàn cục kém khi nền loang sáng nặng.
\end{itemize}

\section{Hướng phát triển}
\label{muc:huongphattrien}
\begin{enumerate}
\item Thu thêm ảnh thật từ nhiều người viết/vẽ khác (hiện chỉ có một người, dù đã đủ số lượng tối thiểu), đánh giá bằng \texttt{python3 danh\_gia.py -{}-bai-toan ca\_hai}, rồi khảo sát nguyên nhân chữ số thật kém hơn MNIST (kiểu viết, cách căn chữ, bước giữ thành phần liên thông lớn nhất).
\item Căn chữ số theo trọng tâm thay vì tâm hộp bao để đường qua tiền xử lý khớp ảnh MNIST hơn (hiện thấp hơn khoảng $1{,}4$ điểm phần trăm trên cùng ảnh).
\item Phân ngưỡng thích nghi cục bộ cho ảnh chụp có ánh sáng không đều.
\item Thêm lớp từ chối ``không phải'', hiệu chuẩn softmax, huấn luyện nhiều hạt để có khoảng tin cậy.
\item Mở rộng tập hình (hình thang, ngũ giác, ngôi sao) và dữ liệu chữ số viết tay của người Việt.
\end{enumerate}


\bibliographystyle{plain}
\bibliography{tai_lieu_tham_khao}

\end{document}
